\documentclass[11pt, oneside]{article} 
\usepackage{geometry}
\geometry{letterpaper} 
\usepackage{graphicx}
	
\usepackage{amssymb}
\usepackage{amsmath}
\usepackage{parskip}
\usepackage{color}
\usepackage{hyperref}

\graphicspath{{/Users/telliott/Dropbox/Github-Math/figures/}}
% \begin{center} \includegraphics [scale=0.4] {gauss3.png} \end{center}

\title{Euclid's Elements}
\date{}

\begin{document}
\maketitle
\Large

%[my-super-duper-separator]

\subsection*{incircle}

\label{sec:incenter}

For any triangle it is possible to draw a circle, called its incircle, which is tangent to all three sides.  When stated in this way, it does not seem obvious how to actually do it.

\begin{center} \includegraphics [scale=0.5] {incircle3.png} \end{center}

\emph{Proof}.

Let $AI$ be the bisector of $A$ and $BI$ be the bisector of $B$, meeting at $I$.

Recall that if we construct the bisector of an angle, say $AI$ bisecting $A$, then every point on the bisector is equidistant from the sides of the original angle.

So one can draw two right triangles with sides $r$ and $AI$, which are congruent by HL, explaining the labels $x$.

The point $I$ is equidistant from side $c$ opposite $C$ and side $b$ opposite $B$.

But the same thing can be done for $BI$ bisecting $B$.

Then draw the circle on center $I$ with radius $r$, and we're done.  This circle is the incircle for $\triangle ABC$.

$r$ is a radius and at the point where it touches the side, perpendicular.  Thus the three sides are each tangent to the incircle.

$\square$

We can obtain two simple results.

First, the area of the triangle is the sum of two copies of each of the three types of right triangle:  $\mathcal{A} = rx + ry + rz$.  But the perimeter of the triangle is twice $x + y + z$, the semi-perimeter $s = x + y + z$ and then

\[ \mathcal{A} = rs \]

Also, we have that 
\[ s = x + y + z \] 
but $y + z = a$ ($a$ is the side opposite the vertex $A$) so:
\[ s = x + a \]
\[ x = s - a \]

The tangent adjacent to vertex $A$ has the length $s - a$.

\subsection*{Pythagoras incircle proof}

\label{sec:PProof_incircle}

Dunham gives this as a problem.

Let right $\triangle ABC$ have sides opposite $a$, $b$ and $c$, as usual.

Let the parts of the sides be $x$ and $y$.  Note that since this is a right triangle, what would be $z$ becomes $r$, the radius of the incircle.

Also, $c = x + y$ is the hypotenuse.

Using the previous result
\[ r = s - c = \frac{a + b - c}{2} \]

\begin{center} \includegraphics [scale=0.35] {pyth25.png} \end{center}

The area of the triangle can be computed as the sum of three quadrilaterals
\[ K = r^2 + rx + ry \]
\[ = r(r + c) \]

Recall that $c = x + y$ and $r = z$ so $r + c = s$, the semiperimeter, or simply remember the result
\[ \mathcal{A} = rs \]

In any event, we have that $r = s - c$.  Substitute for $r$ from above
\[ K = (s - c)s = \frac{a + b - c}{2}  \cdot \frac{a + b + c}{2}  \]

We also have that
\[  K = \frac{ab}{2} \]
so
\[ \frac{ab}{2} = \frac{a + b - c}{2} \cdot \frac{a + b + c}{2}  \]
\[ 2ab = (a + b - c)(a + b + c)  \]
\[ = (a + b)^2 - c^2 \]

Expanding and canceling we get
\[ a^2 + b^2 - c^2 = 0 \]

Rearrange to obtain the Pythagorean relationship.

$\square$

\subsection*{extended angle bisector}

When the angle bisectors are extended to the circumcircle, they form isosceles triangles.

\emph{Proof}.

\begin{center} \includegraphics [scale=0.35] {extended_bisector.png} \end{center}

In the figure above, the incircle is drawn on center $I$ and the angle bisectors as well.  Extend one of them to meet the circumcircle at $D$.

The equal angles produced by bisection are marked with dots.  $\angle BDI = A$ by equal angles on equal arcs.

In $\triangle DBI$, $\angle DBI$ is composed of two angles:  one has a red dot because it lies on an arc equal to that for the half-angle at $C$, often called $\gamma$.

The other one has a blue dot because it is the half-angle of $B$, often called $\beta$.

By sum of angles then, $\angle DIB = \angle DBI = \beta + \gamma$.

It follows that $\triangle DBI$ is isosceles, with $BD = ID$.

A similar result is obtained for $\triangle DAI$

$\square$



\end{document}